%!TEX root = ../thesis.tex

%jama %hakon -- joint introduction chapter

\chapter{Introduction}
\label{ch:intro}

\section{Motivation and Context}
\label{sec:motivation}

The COVID-19 pandemic showed that infectious disease dynamics operate at more than one scale at the same time \citep{dong2022,mathieu2020}.
At the population level, the average transmission rate determines whether an epidemic grows or shrinks.
At the individual level, contact patterns, mobility, vaccination status, and personal risk decide who gets infected, hospitalised, or dies \citep{zhang2021}.
A control strategy designed using only a population-level model can be correct on average.
It may still miss local outbreaks, regional hospital overload, or the effect of high-mobility individuals on cross-city spread.

Control theory has been applied to epidemic management in several recent works.
Model Predictive Control (MPC) in particular has been used to compute intervention policies that respect hospital capacity while limiting societal impact.
For population-level ordinary differential equation (ODE) models, MPC schemes can drive the epidemic to a disease-free equilibrium with formal stability guarantees \citep{esterhuizen2024}.
Whether the same MPC structure also works when the actual plant is a stochastic, agent-level system has received less attention.
This is the gap addressed in this thesis.

This thesis is a joint bachelor project with two authors.
The macro-level part, by H\aa{}kon Rullestad H\aa{}varstein, is a multi-city SIRVHDB compartmental model in MATLAB/Simulink.
The macro MPC operates in closed loop with the nonlinear ODE plant inside Simulink.
The micro-level part, by Mohamed Ahmed Jama, is a stochastic agent-based model (ABM) of the same two-city population.

For the cross-scale evaluation, a second MPC controller is implemented in MATLAB.
Its internal prediction model is a linearised version of the SIRVHDB ODE.
The linearisation uses parameters fitted by a genetic algorithm (GA) that aligns the ODE response with ABM trajectories.
At each daily step, this controller reads the ABM's combined hospital load, computes the per-city intervention signals $u_1(t)$ and $u_2(t)$, and the resulting transmission rates $\beta_1(t)$ and $\beta_2(t)$ are applied to the ABM in the next step.
The ABM is therefore the plant that the MPC controls in closed loop.
The mismatch between the linearised ODE used for prediction and the stochastic ABM used as the actual plant is part of what this thesis evaluates.

\section{Problem Statement}
\label{sec:problem}

This thesis addresses two problems.

The first problem is macro-level epidemic control.
Given a multi-city SIRVHDB compartmental model with COVID-19-based parameters, can an MPC controller keep combined hospital load below a chosen capacity threshold, suppress secondary infection waves, and keep the intervention signal within realistic operational limits?

The second problem is cross-scale validation.
The MPC controller uses a linearised ODE as its internal prediction model, and the agent-based model is the actual plant being controlled.
Can this controller still keep hospital load below the capacity threshold and suppress epidemic dynamics, given the mismatch between its linear predictor and the stochastic, individual-level dynamics of the plant?

Both problems depend on a supporting task: parameter calibration.
The macro ODE and the ABM represent epidemic dynamics in different ways, so shared nominal parameters do not guarantee consistent behaviour.
A systematic GA-based calibration procedure is needed to align the two models before the linearised plant can be used as the prediction model of the cross-scale MPC controller.

\section{Scope and Delimitations}
\label{sec:scope}

The macro model is a two-city SIRVHDB system in MATLAB/Simulink.
The macro MPC controller is built with the MATLAB MPC Toolbox and runs in closed loop with the nonlinear ODE plant inside Simulink.
All macro simulations use a single 500-day horizon.
The first 150 days of each run are reported as the early-phase response.
The full 500 days are reported as the long-horizon response.

The micro model is implemented in MATLAB.
It runs for $1{,}000$ agents (500 per city) over a 150-day horizon.
The cross-scale MPC controller is built separately from the Simulink macro MPC.
It runs in closed loop with the ABM and uses a GA-fit linearisation of the SIRVHDB ODE as its internal prediction model.

Each agent is assigned to a home city at initialisation.
Inter-city movement is modelled as temporary commuting, not permanent migration.
The home-city populations $N_1 = N_2 = 500$ remain constant throughout the simulation.
Deceased agents remain in their home city's bookkeeping.

Both models use parameter values drawn from the COVID-19 literature (transmission rate, recovery rate, hospitalisation rate, vaccine efficacy).
Neither model is fitted to a specific real-world outbreak.
The goal is to study control behaviour in an epidemiologically plausible regime, not to reproduce historical case data.

\section{Contributions}
\label{sec:contributions}

The contributions of this thesis are:
\begin{itemize}
  \item A two-city SIRVHDB compartmental epidemic model in MATLAB/Simulink, with waning immunity, seasonal transmission forcing, and inter-city coupling through a travel matrix and commuting-based mixing.
  \item A macro-level MPC controller that operates in closed loop with the Simulink ODE plant, regulating combined hospital load below a capacity threshold using per-city transmission reduction as the control input, supported by a Kalman filter state estimator.
  \item A spatially explicit stochastic agent-based model of the same two-city population, with seven disease states, proximity-based infection, vaccination with breakthrough infections, and inter-city commuting.
  \item Implementation of inter-city commuting as a probabilistic transfer between two abstract regions, generating cross-city epidemic coupling without modelling real geographic coordinates or transport networks.
  \item A genetic-algorithm (GA) calibration procedure that fits ODE parameters to ABM hospital-load trajectories across several test scenarios, producing a calibrated ODE that is then linearised to serve as the prediction model of the cross-scale MPC controller.
  \item A closed-loop integration in which the cross-scale MPC controller, using the GA-fit linearised plant as its internal model, controls the ABM directly. The ABM's hospital load is fed back to the controller at each step, and the controller's per-city transmission reduction signals are applied to the ABM in the next step.
  \item A simulation-based evaluation of MPC performance at both modelling levels, covering macro-level wave suppression and capacity adherence over a 500-day horizon, and cross-scale confirmation that an MPC built from a GA-fit linearised ODE can regulate the stochastic ABM in closed loop despite the model mismatch.
\end{itemize}

\section{Thesis Outline}
\label{sec:outline}

The remainder of this thesis is organised as follows.

Chapter~\ref{ch:related} reviews relevant literature on compartmental epidemic models, agent-based epidemic models, multi-scale epidemic modelling, and MPC applied to epidemic control.

Chapter~\ref{ch:macro_model} presents the macro-level SIRVHDB model, the multi-city interaction mechanisms (travel matrix and commuting-based mixing), the macro MPC formulation, and the Kalman filter state estimator.

Chapter~\ref{ch:macro_results} presents the macro-level simulation results over a 500-day horizon, covering baseline epidemic dynamics, city-to-city interaction scenarios, waning immunity effects, and macro MPC controller performance in both the early phase and the long-horizon regime.
This chapter answers Problem~1.

Chapter~\ref{ch:integration} describes the macro--micro integration architecture: the GA calibration of ODE parameters to ABM trajectories, the linearisation step that produces the prediction model used by the cross-scale MPC controller, and the closed-loop coupling between this controller and the ABM.

Chapter~\ref{ch:abm} presents the design and implementation of the agent-based model, including agent states, disease transmission and progression, the two-city spatial structure, and the aggregation interface used by the cross-scale controller.

Chapter~\ref{ch:micro_results} presents the micro-level simulation results, covering the baseline uncontrolled epidemic, the effect of varying transmission rates, the two-city comparison, and the GA calibration results.

Chapter~\ref{ch:closed_loop} presents the closed-loop simulation results in which the cross-scale MPC controller drives the ABM.
Controlled and uncontrolled scenarios are compared in terms of infection dynamics, hospital load, mortality, and capacity adherence.
This chapter answers Problem~2.

Chapter~\ref{ch:discussion} interprets the results from both model levels, discusses the value of the multi-scale approach, and notes the limitations of the current framework.

Chapter~\ref{ch:conclusion} concludes the thesis with a summary of findings and directions for future work.

Appendix~\ref{apx:instructions} gives instructions for compiling and running the macro Simulink model and the cross-scale MATLAB simulation.